% Questões da AV2 - Containerização e CI da aplicação Copa Resultados
% Os gabaritos das questões são lidos diretamente de Test/snippets/,
% que contém cópias exatas dos arquivos reais do projeto App/.

% 1 - App/api/Dockerfile
\begin{questao}[0,5 ponto]
Escreva o conteúdo completo do arquivo \texttt{App/api/Dockerfile}. O Dockerfile
deve ser \emph{multi-stage}, com duas etapas:

\begin{itemize}
  \item \textbf{\texttt{build}}: instala todas as dependências e compila o
        TypeScript;
  \item \textbf{\texttt{production}}: instala somente as dependências de
        produção e executa a API já compilada.
\end{itemize}

\resposta[22]{Conteúdo exato de \texttt{App/api/Dockerfile}:
\lstinputlisting{snippets/api-Dockerfile.txt}}
\end{questao}

% 2 - App/docker-compose.yml
\begin{questao}[0,5 ponto]
Escreva o conteúdo completo do arquivo \texttt{App/docker-compose.yml},
orquestrando os serviços \texttt{mongodb}, \texttt{api} e \texttt{web}.
Considere que:

\begin{itemize}
  \item as portas expostas no host e as demais configurações devem usar as
        variáveis de ambiente do arquivo \texttt{.env} da raiz
        (apresentadas na contextualização);
  \item os dados do MongoDB devem persistir em um volume nomeado;
  \item o serviço \texttt{api} depende do \texttt{mongodb}, e o serviço
        \texttt{web} depende da \texttt{api};
%   \item o serviço \texttt{web} deve repassar \texttt{VITE\_API\_URL} como
%         \emph{build arg}.
\end{itemize}

\resposta[54]{Conteúdo exato de \texttt{App/docker-compose.yml}:
\lstinputlisting{snippets/docker-compose.txt}}
\end{questao}

% App/.env
% \begin{questao}[0,2 ponto]
% Com base nas variáveis de ambiente apresentadas na contextualização, escreva o
% conteúdo completo do arquivo \texttt{App/.env}, utilizado pelo Docker
% Compose.

% \resposta[10]{Conteúdo exato de \texttt{App/.env}:
% \lstinputlisting{snippets/env-dev.txt}}
% \end{questao}

% 3 - App/.github/workflows/api-test-bdd.yml
\begin{questao}[0,7 ponto]
Escreva o conteúdo completo do arquivo
\texttt{App/.github/workflows/api-test-bdd.yml}, o workflow de
\emph{GitHub Actions} que executa os testes BDD da API a cada \emph{push} no
repositório. Considere que:

\begin{itemize}
  \item o workflow chama-se \texttt{API - Testes BDD} e deve disparar em
        qualquer \emph{push};
  \item há um único \emph{job} chamado \texttt{test-bdd} (nome de exibição
        \texttt{API - testes BDD}) que roda em \texttt{ubuntu-latest};
  \item o \emph{job} usa um \emph{service container} de banco de dados com a
        imagem \texttt{mongo:7}, publicando a porta \texttt{27017:27017};
  \item por padrão, todos os comandos \texttt{run} devem executar no diretório
        \texttt{App/api} (\texttt{working-directory});
  \item o \emph{job} define as variáveis de ambiente
        \texttt{MONGODB\_URI=mongodb://localhost:27017/cupmatches},
        \texttt{PORT=3000} e \texttt{API\_URL=http://localhost:3000};
  \item os passos, nesta ordem, são: 
      \begin{itemize}
            \item \emph{checkout} do código (\texttt{actions/checkout@v4});
            \item configurar o Node.js 20 (\texttt{actions/setup-node@v4});
            \item instalar as dependências;
            \item compilar o TypeScript;
            \item popular o banco com o \emph{seed};
            \item subir a API em segundo plano e aguardar 5 segundos; e
\begin{lstlisting}
run: |
  npm start &
  sleep 5
\end{lstlisting}
            \item executar os testes BDD.
      \end{itemize}
\end{itemize}


% Os \emph{scripts} npm disponíveis são: \texttt{npm install}, \texttt{npm run
% build}, \texttt{npm run seed}, \texttt{npm start} e \texttt{npm run test:bdd}.

\resposta[54]{Conteúdo exato de
\texttt{App/.github/workflows/api-test-bdd.yml}:
\lstinputlisting{snippets/github-actions.txt}}
\end{questao}

% 5 - Sequência de comandos
\begin{questao}[0,3 ponto]
Considere o seguinte cenário: o Docker está iniciado; todos os pacotes do
projeto estão instalados; o link do repositório remoto já foi adicionado ao
projeto; e todas as mudanças estão comitadas. Escreva, na ordem correta, a
sequência de comandos para \textbf{subir os containers} e, em seguida,
\textbf{atualizar o repositório remoto}.

\resposta[10]{Sequência de comandos:
\lstinputlisting{snippets/comandos.txt}
Observação: o \emph{push} dispara automaticamente o workflow
\texttt{api-test-bdd.yml} no GitHub Actions, que executa os testes BDD da API.}
\end{questao}
